\documentclass[12pt]{article}
\usepackage[margin=1in]{geometry} % narrower margins
\usepackage[UTF8,fontset=fandol]{ctex}  
\renewcommand{\figurename}{Figure}
\renewcommand{\tablename}{Table}
\usepackage{amsmath, amssymb, graphicx, setspace, natbib}
\usepackage{booktabs,tabularx}
\usepackage{multirow}
\usepackage{comment}
\usepackage{enumitem}
\usepackage{array}

\title{Tidal Lanes}
\author{Xiaoruo Feng}
\date{February 2026}


\begin{document}
	\maketitle
	
	\section{Introduction}
	\begin{itemize}
		\item Motivation: asymmetric traffic flows (tidal lanes) versus symmetric assumptions in standard spatial models.
		\item Contribution: extend Allen \& Arkolakis (2022) to incorporate directional, grid-level road networks.
		\item Data advantage: high-frequency road segment speed data (10-minute interval) combined with commuting-flow matrices (Beijing and Shenzhen).
	\end{itemize}
	
	\section{Data}
	\subsection{Commuting-Flow Data}
	
	\begin{itemize}
		\item \textbf{Beijing:} 2022 100-meter gird OD data
		\begin{itemize}
			\item File path: \par  \texttt{/Users/fxr/Dropbox/1Data/Baidu/CommutingBeijing/Data\_2021\_2022\_FromBICP}.
			\item Observations: 6,909,608 (O-D pairs). Among them, there are 355,190 unique origin grid IDs, 394,500 unique destination grid IDs, and a total of 453,391 unique grids.
			\item Size: 100-meter grids (WGS84 coordinates; in practice the resolution is approximately 75 meters per grid cell).
			\item Infor: Commuters are disaggregated by demographic and socioeconomic characteristics, including gender, age, education, income, consumption patterns, and occupation categories.
		\end{itemize}

		\item \textbf{Shenzhen:} 2024 100-meter grid OD data
		\begin{itemize}
			\item File path:  \texttt{/Users/fxr/Dropbox/1Data/Baidu/CommutingShenzhen/Data}
			\item Observations: 3,970,569 (O-D pairs). Among them, there are 94,670 unique origin grid IDs, 105,322 unique destination grid IDs, and a total of 110,339 unique grids.
			\item Size: 100-meter grids
			\item Infor: Commuters are disaggregated by demographic and socioeconomic characteristics, including gender, age, education, income, consumption patterns, and occupation categories.
		\end{itemize}
				
		\item \textbf{Limitation:} Both commuting matrices do not distinguish between different transport modes.
	\end{itemize}

	\begin{figure}[h!]
		\centering
		\begin{minipage}[t]{0.49\textwidth}
			\centering
			\includegraphics[width=\textwidth]{../../L3_figs/grids_work_pop_beijing_fixed.pdf}
		\end{minipage}
		\hfill
		\begin{minipage}[t]{0.49\textwidth}
			\centering
			\includegraphics[width=\textwidth]{../../L3_figs/grids_resi_pop_beijing_fixed.pdf}
		\end{minipage}
		\caption{Working and Residential Distribution in Beijing}
		\label{fig:Beijing_commuting}
	\end{figure}
		
	\subsection{Road Segment Speed Data}
	\begin{itemize}
		\item Beijing: segment-level speed, 10-minute frequency, with attributes \texttt{cityname}, \texttt{roadseg\_id}, \texttt{roadname}, \texttt{roadtype}, \texttt{semantic}, \texttt{location}, \texttt{geometry}.
		\item Shenzhen: same structure.
	\end{itemize}
	

	\begin{table}[h!]
		\centering
		\caption{Example record from the raw road-segment dataset}
		\label{tab:raw_example}
		\small
		\setlength{\tabcolsep}{6pt}
		\begin{tabularx}{\textwidth}{@{}l X l >{\raggedright\arraybackslash}p{0.55\textwidth}@{}}
			\toprule
			\textbf{Variable} & \textbf{Value} & \textbf{Variable} & \textbf{Value} \\
			\midrule
			cityname &
			北京市 &
			roadseg\_id &
			\texttt{RS\_NB1520263513\_NB1530878951\_1} \\
			roadname &
			\texttt{X017} &
			semantic &
			自 S320 出入口到 X017 路口，东向西 \\
			roadtype &
			4 &
			geometry &
			\texttt{LINESTRING(115.70157 39.57762, 115.70108 39.5\ldots)} \\
			location &
			0.043 &
			& \\			
			\bottomrule
		\end{tabularx}
	\end{table}

	\begin{figure}[h!]
		\centering
		\includegraphics[width=0.9\textwidth]{../../L3_figs/speed_eg.png}
		\label{fig:speed_eg}
	\end{figure}

	\clearpage

%%%%%%%%%%%%%%%%%%%%%%%%%%%%%%%%%%%%%%%%%	

	\section{Road Network Construction and Centerline Matching}
	
	The raw road map is rich enough for measurement but too irregular to be used directly as the directed network in the main analysis. Divided carriageways are frequently offset, ramps and frontage roads generate local duplication, and grade-separated interchanges create stacked geometries that are difficult to reconcile with a one-segment-per-corridor representation. We therefore map the raw road system into a direction-aware centerline network and retain an explicit crosswalk from raw segments to centerlines. The appendix has two purposes. The first is to describe the measurement object. The second is to show that the resulting match is tight enough for the speed aggregation used below.
	
	\subsection{Raw Road Network and Direction Information}
	
	The construction starts from the raw Beijing road polylines. After projection to a planar coordinate system, the raw network serves as the geometric source for all subsequent measurement. Each segment also retains the directional information carried by the original map attributes, which is later used when assigning observations to the directed centerline network.
	
	Figure~\ref{fig:raw-centerline-compare} illustrates the role of this transformation. The raw map contains multiple parallel traces for the same corridor and local geometric noise around ramps and intersections. The centerline representation collapses those local duplications into a cleaner object that is better suited for directional measurement.
	
	\begin{figure}[h!]
		\centering
		\begin{minipage}[t]{0.49\textwidth}
			\centering
			\includegraphics[width=\textwidth]{../../../00archive/exhibit/map_raw_roads_15km.png}
		\end{minipage}
		\hfill
		\begin{minipage}[t]{0.49\textwidth}
			\centering
			\includegraphics[width=\textwidth]{../../../00archive/exhibit/map_centerline_15km.png}
		\end{minipage}
		\caption{Raw roads and derived centerlines within 15 km of Tiananmen Square. The figure shows why a centerline abstraction is useful: the raw network contains local duplication and irregular geometry, while the centerline network recovers the corridor-level structure needed for directional aggregation.}
		\label{fig:raw-centerline-compare}
	\end{figure}
	
	\subsection{Centerline Extraction via Skeletonization}
	
	The retained construction uses a 50-meter road-surface buffer, a 5-meter raster for skeletonization, and a 50-meter minimum length threshold for spurious branches. Those choices are conservative. The buffer is wide enough to merge parallel traces that belong to the same physical corridor, the 5-meter raster preserves curvature better than a coarser grid, and the short-branch pruning removes small artifacts that arise near complex intersections.
	
	\subsection{Directed Centerline Network Construction}
	
	Once the undirected centerline is extracted, each centerline segment is duplicated into two directions. The resulting network is symmetric by construction, which is useful for directional speed aggregation because each physical corridor admits one link in each direction. This representation does not attempt to preserve every lane-level feature in the raw map. It recovers the corridor-level object on which directional travel conditions are measured.
	
	\subsection{Selective Segmentation of Raw Road Segments}
	
	Before matching, raw edges are split only when local geometry indicates that a single segment spans more than one centerline element. The rule is deliberately conservative. In the retained construction, the segment count rises from 24,763 raw edges to 24,767 split segments. The procedure therefore is not relying on aggressive segmentation to improve the match rate.
	
	\subsection{Matching Raw Sub-Segments to Directed Centerlines}
	
	Skeletonization works well for the bulk of the network, but the failure cases are systematic rather than random. They appear most often on long divided expressways, curved frontage roads, ramps, and grade-separated interchanges where the two directions do not sit symmetrically around a stable medial axis. In these settings, a purely mechanical skeleton can attach the centerline to the wrong carriageway, cut a corridor into implausibly short pieces, or bridge across upper and lower structures that should remain distinct. We therefore supplement the default skeleton with a manually curated set of corridors that are matched through an exact raw-to-centerline crosswalk.
	
	\subsection{Outputs and Descriptive Statistics}
	
	Five related network objects enter the measurement pipeline. The first two are the raw road network and its selectively split counterpart. The next two are the undirected and directed centerline systems before any baseline restriction. The fifth, which is the network used in the speed and lane stages, is the baseline-kept directed centerline network. The baseline rule is geometric: a directed centerline is retained when it is a valid \texttt{LineString} and its parent centerline is at least 50 meters long. Table~\ref{tab:network-summary} and Figure~\ref{fig:length-cdf-5networks} summarize how these objects differ.
	
	\begin{table}[h!]
		\centering
		\caption{Network summary statistics across alternative representations.}
		\label{tab:network-summary}
		\begin{tabular}{l r r r r}
			\toprule
			Network & $N$ & Mean length (m) & Median length (m) & Total length (km) \\
			\midrule
			raw & 23,799 & 792 & 428 & 18,850.8 \\
			raw\_split & 23,799 & 792 & 428 & 18,850.8 \\
			undirected centerline (all) & 22,876 & 429 & 7 & 9,812.4 \\
			directed centerline (all) & 45,752 & 429 & 7 & 19,624.9 \\
			directed centerline (baseline) & 11,890 & 1,631 & 708 & 19,391.3 \\
			\bottomrule
		\end{tabular} \\
		\vspace{0.3em}
		{\footnotesize\raggedright
			Note: the baseline network keeps directed centerlines whose geometry is valid and whose parent undirected centerline is at least 50 meters long. The restriction removes short skeleton branches and leaves the corridor backbone used in downstream speed aggregation.
		}
	\end{table}
	
	The contrast between the unrestricted and baseline directed networks is informative rather than cosmetic. Before the length screen, the centerline system contains many short skeleton branches, which is why the unrestricted directed network has a median length of only 7 meters. After the baseline restriction, the median rises to 708 meters and the total retained length barely changes. The filter therefore removes spurious local branches, not economically meaningful corridors.
	
	\begin{figure}[h!]
		\centering
		\includegraphics[width=0.74\textwidth]{../../../data_work/outputs/manual_centerline_rules_v11/figures/paper_appendix/fig_length_cdf_5networks_v11.png}
		\caption{Length CDF for the raw, raw\_split, undirected centerline, directed centerline, and baseline directed centerline networks. The baseline restriction trims short branches while preserving almost all corridor length.}
		\label{fig:length-cdf-5networks}
	\end{figure}
	
	The key empirical point is that the manually handled corridors are not a large share of the network by count, but they are a large share by length. Table~\ref{tab:manual-corridor-summary} shows that the manually handled edges account for only 963 raw edges, or 3.9\% of all raw edges, yet they represent 3,846 km of roadway, which is 20.4\% of total split-segment length. The typical manually handled edge is much longer than the rest of the network. Its mean length is about 4.0 km, compared with 0.63 km for the remaining edges. This is exactly what one would expect if manual intervention is concentrated on long, complex corridors rather than on local fragments.
	
	\begin{table}[h!]
		\centering
		\caption{Why manual corridors matter.}
		\label{tab:manual-corridor-summary}
		\begin{tabularx}{\textwidth}{>{\raggedright\arraybackslash}X r r r r r}
			\toprule
			Group & Count & Edge share & Length (km) & Mean (m) & Median (m) \\
			\midrule
			Manual exact-match corridors & 963 & 3.9\% & 3,846.0 & 3,994 & 1,252 \\
			Remaining raw edges & 23,800 & 96.1\% & 15,026.7 & 631 & 400 \\
			\bottomrule
		\end{tabularx} \\
		\vspace{0.3em}
		{\footnotesize\raggedright
			Note: ``Manual exact-match corridors'' refers to raw edges assigned through the manual-corridor crosswalk. The table does not mean that one fifth of the network is hand-drawn. It means that about one fifth of split-segment length is attached through the manual exact-match channel.
		}
	\end{table}
	
	Figure~\ref{fig:manual-unmatched-diagnostics} complements Table~\ref{tab:manual-corridor-summary}. Panel A highlights the manually handled corridors. They are concentrated on major, elongated corridors rather than on short local streets. Panel B plots the remaining unmatched edges longer than 500 meters. The residual is sparse and concentrated in a small set of difficult roads, which is consistent with the view that the remaining error is a narrow long tail rather than a citywide mismatch problem.
	
	\begin{figure}[h!]
		\centering
		\includegraphics[width=0.92\textwidth]{../../../data_work/outputs/manual_centerline_rules_v11/figures/paper_appendix/manual_vs_unmatched_corridors.png}
		\caption{Manual-corridor coverage and remaining long unmatched corridors. The left panel highlights corridors assigned through the manual exact-match channel. The right panel plots the remaining unmatched raw edges longer than 500 meters.}
		\label{fig:manual-unmatched-diagnostics}
	\end{figure}
	
	\subsection{Match Assessment: Length Coverage and Matching Statistics}
	
	Table~\ref{tab:match-summary-v11} reports the quality of the retained raw-to-centerline crosswalk. On the baseline-kept split network, 23,662 of 23,799 split segments are matched, which yields a 99.42\% match rate. On the raw-edge network, 23,672 of 24,763 edges are matched. The count-based unmatched share is therefore 4.41\% on raw edges, but the length-based margin is much smaller: unmatched edges account for only 239.9 km out of 18,872.7 km of raw-network length, or 1.27\%.
	
	\begin{table}[h!]
		\centering
		\caption{Matching quality of the raw-to-centerline crosswalk.}
		\label{tab:match-summary-v11}
		\begin{tabular}{l r r r r}
			\toprule
			Measure & Total & Matched & Unmatched & Rate / share \\
			\midrule
			Baseline-kept split segments & 23,799 & 23,662 & 137 & 99.42\% matched \\
			All raw edges & 24,763 & 23,672 & 1,091 & 95.59\% matched \\
			Raw-network length (km) & 18,872.7 & 18,632.8 & 239.9 & 1.27\% unmatched \\
			\bottomrule
		\end{tabular}
	\end{table}
	
	We next reverse the lens and ask whether a centerline corridor receives support from raw observations. Coverage is not the main quality metric, because the centerline network is intentionally more complete than the matched speed network. It is nonetheless useful for understanding why the remaining lane and speed gaps are concentrated in short tails rather than in the arterial backbone.
	
	\begin{table}[h!]
		\centering
		\caption{Coverage of directed centerlines by raw sub-segments.}
		\label{tab:directed-coverage}
		\begin{tabular}{l r r r r}
			\toprule
			Direction & $N$ & Covered & Uncovered & Coverage rate \\
			\midrule
			ALL & 11,890 & 8,393 & 3,497 & 0.7059 \\
			AB  & 5,945 & 4,229 & 1,716 & 0.7114 \\
			BA  & 5,945 & 4,164 & 1,781 & 0.7004 \\
			\bottomrule
		\end{tabular}
	\end{table}
	
	\begin{table}[h!]
		\centering
		\caption{Coverage status of undirected centerlines.}
		\label{tab:undirected-coverage}
		\begin{tabular}{l r r}
			\toprule
			Status & $N$ & Share \\
			\midrule
			Both AB and BA & 3,697 & 0.6219 \\
			AB only or BA only & 999 & 0.1680 \\
			None & 1,249 & 0.2101 \\
			\bottomrule
		\end{tabular}
	\end{table}
	
	The directional-coverage pattern sharpens this point. On the baseline centerline backbone, 62.2\% of undirected corridors are matched in both directions, 16.8\% are matched in only one direction, and 21.0\% receive no directional support. The length perspective is far more favorable. Corridors matched in both directions account for 87.0\% of baseline centerline length, one-direction corridors account for another 9.9\%, and corridors unmatched in both directions account for only 3.0\% of baseline length. From the raw-edge perspective, 90.0\% of edges, representing 91.9\% of raw length, fall on corridors whose centerlines are supported in both directions.
	
	The residual unmatched network is highly uneven. Most unmatched edges are short fragments. Among the 1,091 unmatched raw edges, 954 are shorter than 50 meters. These account for 87.4\% of unmatched edges by count but only 11.9\% of unmatched length. The opposite is true for the long tail: only 134 unmatched edges are at least 500 meters long, yet they account for 87.7\% of unmatched length. The remaining error is therefore concentrated in a narrow set of corridor-level cases rather than diffused throughout the street grid.
	
	\begin{figure}[h!]
		\centering
		\begin{minipage}[t]{0.49\textwidth}
			\centering
			\includegraphics[width=\textwidth]{../../../data_work/outputs/manual_centerline_rules_v11/figures/final_match_review/map_final_match_centerline_coverage.png}
		\end{minipage}
		\hfill
		\begin{minipage}[t]{0.49\textwidth}
			\centering
			\includegraphics[width=\textwidth]{../../../data_work/outputs/manual_centerline_rules_v11/figures/final_match_review/map_final_match_raw_edges_4status.png}
		\end{minipage}
		\caption{Directional match status on baseline centerlines and raw edges. The left panel classifies baseline centerlines by whether AB only, BA only, both directions, or neither direction is supported by matched raw segments. The right panel assigns raw edges to the same corridor-level classification.}
		\label{fig:matched-unmatched-map}
	\end{figure}
	
	Taken together, the diagnostics support a simple conclusion. The centerline abstraction is necessary because the raw network is too irregular for direct directional aggregation. Manual intervention is necessary because a small set of long, geometrically complex corridors systematically defeats a generic skeleton. Once those corridors are handled directly, the remaining mismatch is modest in count and very small in length, which is the relevant margin for the speed measures used in the rest of the paper.
	
	\subsection{Lane-Count Proxy on the Baseline Centerline Network}
	
	The lane proxy is constructed after matching, not from the speed stage. Each matched raw segment contributes its road-class-based lane count to the directed centerline on which it lands, and the aggregation is length-weighted over matched overlap. In the retained run, the road-class field is complete on matched raw segments. Missing lane estimates therefore arise because a baseline centerline receives no raw match, not because the raw attribute is absent.
	
	Table~\ref{tab:lane-summary} summarizes the resulting coverage. On the baseline directed network, 8,393 of 11,890 links receive a lane estimate, which is 70.6\% by count and 92.0\% by length. The mean estimated lane count is 3.13 and the median is 3.32. Opposite-direction estimates are close on average, with a mean lane ratio of 1.04; among links with valid estimates, 86.2\% have an opposite-direction ratio below 1.5. Figure~\ref{fig:lane-coverage} maps the baseline coverage pattern. The missing links are visible, but they are concentrated in short peripheral pieces rather than in the main arterial backbone.
	
	\begin{table}[h!]
		\centering
		\caption{Lane-count proxy on the baseline directed centerline network.}
		\label{tab:lane-summary}
		\begin{tabular}{l r}
			\toprule
			Statistic & Value \\
			\midrule
			Baseline directed centerlines & 11,890 \\
			With lane estimate & 8,393 \\
			Count coverage & 70.6\% \\
			Length coverage & 92.0\% \\
			Mean lane estimate & 3.13 \\
			Median lane estimate & 3.32 \\
			Mean opposite-direction lane ratio & 1.04 \\
			Symmetric share among valid estimates & 86.2\% \\
			\bottomrule
		\end{tabular}
	\end{table}
	
	\begin{figure}[h!]
		\centering
		\includegraphics[width=0.76\textwidth]{../../../data_work/outputs/manual_centerline_rules_v11/figures/paper_appendix/fig_lane_coverage_baseline_v11.png}
		\caption{Baseline directed centerlines with and without lane estimates. Blue links receive a lane-count proxy from matched raw road classes; orange links do not.}
		\label{fig:lane-coverage}
	\end{figure}

		\clearpage

%%%%%%%%%%%%%%%%%%%%%%%%%%%%%%%%%%%%%%%%%	
	\section{Aggregating Raw-Segment Speeds to Directed Centerlines}
	
	Raw speed observations are attached to directed centerlines through the matched crosswalk developed in Section~3. The aggregation is carried out at the weekday-status by hour level. In the retained run, 23.29 million matched raw speed observations are mapped into 394,468 directed-centerline-hour cells spanning 8,349 unique directed centerlines. The purpose of this stage is to construct directional travel speeds on a corridor object that is clean enough for later asymmetry and grid aggregation.
	
	\paragraph{Time-weighted aggregation.}
	
	We compute the average speed on each directed centerline using a time-weighted scheme that treats the total distance as $\sum_{r \in \mathcal{R}(c)} \ell_r$ and the total travel time as $\sum_{r \in \mathcal{R}(c)} \ell_r / v_r$, where $\mathcal{R}(c)$ denotes the set of matched raw segments for centerline $c$. The average speed is then:
	
	\begin{align}
		\bar{v}_c^{\,\text{time}}
		&= \frac{\sum_{r \in \mathcal{R}(c)} \ell_r}{\sum_{r \in \mathcal{R}(c)} \ell_r / v_r}.
		\label{eq:centerline_speed}
	\end{align}
	
	This method divides total distance by total time to obtain the average speed, which is the natural aggregation when considering travel time along a centerline. The computation is performed on the matched and quality-filtered set $\mathcal{R}(c)$.
	
	\paragraph{Distributional and diurnal profiles.}
	
	The aggregate distribution remains plausible after the crosswalk and harmonic aggregation. Figure~\ref{fig:centerline-speed-dist} plots the weekday and weekend distributions of directed centerline speeds. The two distributions are close, with a modest right shift on weekends, which is what one would expect in a city where the strongest directional congestion is tied to commuting peaks rather than to all-day network collapse.
	
	The within-day profile is sharper. Figure~\ref{fig:centerline-speed-diurnal} plots centerline-length-weighted hourly means. Weekday speeds fall in both commuting peaks, with a deeper trough in the evening. Weekend speeds are flatter throughout the day. The period means line up with the same pattern: the AM mean is 35.19 km/h, the PM mean is 33.33 km/h, and the late-night free-flow window averages 38.46 km/h.
	
	\begin{figure}[h!]
		\centering
		\includegraphics[width=0.68\textwidth]{../../../data_work/outputs/manual_centerline_rules_v11/figures/paper_appendix/fig_centerline_speed_distribution_v11.png}
		\caption{Distribution of directed centerline speeds in the retained run, split by weekday and weekend.}
		\label{fig:centerline-speed-dist}
	\end{figure}
	
	\begin{figure}[h!]
		\centering
		\includegraphics[width=0.82\textwidth]{../../../data_work/outputs/manual_centerline_rules_v11/figures/paper_appendix/fig_centerline_speed_diurnal_v11.png}
		\caption{Directed centerline speed over the day, using centerline-length-weighted hourly means. Shaded bands mark the AM and PM commuting peaks.}
		\label{fig:centerline-speed-diurnal}
	\end{figure}
	
	\subsection{Summary Statistics of Speed}
	
	Table~\ref{tab:stage03-summary} reports the main operational outputs of the speed-aggregation stage. Table~\ref{tab:centerline-speed-stats} then summarizes the resulting directed centerline speeds. The length-weighted mean is above the simple mean in all cases, which implies that longer corridors tend to carry faster traffic. The gap is modest, not dramatic, which is consistent with the view that aggregation is preserving the citywide speed distribution rather than inventing a new one.
	
	\begin{table}[h!]
		\centering
		\caption{Stage03 aggregation summary.}
		\label{tab:stage03-summary}
		\begin{tabular}{l r}
			\toprule
			Statistic & Value \\
			\midrule
			Matched raw speed observations & 23,286,622 \\
			Directed centerline-hour rows & 394,468 \\
			Unique directed centerlines & 8,349 \\
			Overall mean speed (km/h) & 36.30 \\
			Overall median speed (km/h) & 33.33 \\
			AM mean speed (km/h) & 35.19 \\
			PM mean speed (km/h) & 33.33 \\
			22:00--05:00 mean speed (km/h) & 38.46 \\
			\bottomrule
		\end{tabular}
	\end{table}
	
	\begin{table}[h!]
		\centering
		\caption{Summary statistics of directed centerline speeds (km/h).}
		\label{tab:centerline-speed-stats}
		\begin{tabular}{lccccc}
			\toprule
			Group & $N$ & Mean & Weighted Mean & Median & Std. Dev. \\
			\midrule
			All & 394,468 & 36.30 & 40.88 & 33.33 & 12.54 \\
			Weekday & 198,300 & 36.12 & 41.02 & 33.23 & 12.66 \\
			Weekend & 196,168 & 36.47 & 40.75 & 33.43 & 12.42 \\
			\bottomrule
		\end{tabular} \\
		\vspace{0.3em}
		{\footnotesize\raggedright
			Note: the weighted mean uses directed centerline length as weights.
		}
	\end{table}
	
	\clearpage

	%%%%%%%%%%%%%%%%%%%%%%%%%%%%%%%%%%%%%%%%%	
	
	\section{Identifying Asymmetric Segments and Tidal Lanes in Road Network}
	\label{sec:asym_sec}
	This section documents the construction of peak-period directional speed measures on the centerline network and the identification of directional asymmetries and tidal lane candidates. The focus is on transparent data processing and classification rules rather than behavioral interpretation.
	
	\subsection{Time Aggregation and Peak-Period Speed Computation}
	
	Speed observations vary systematically over time. We aggregate raw speed observations by centerline segment, direction, weekday status, and hour of day using a time-weighted harmonic mean (Equation~\eqref{eq:centerline_speed}). The analysis focuses on weekday peak hours only:
	\begin{itemize}
		\item Morning peak (AM): 7:00--9:00
		\item Evening peak (PM): 17:00--19:00
	\end{itemize}
	Within each peak window, observations are aggregated across the two constituent hours to obtain a single peak-period speed for each directed centerline segment, yielding peak-period speeds $\bar{v}_{c,d,p}$ for centerline $c$, direction $d \in \{AB, BA\}$, and peak period $p \in \{AM, PM\}$.
	
	The study area is defined as all directed centerline segments within a 150 km radius of Tiananmen Square in Beijing. After aggregation, valid peak-period observations are available for 4,656 undirected centerlines in the AM peak and 4,659 in the PM peak. Among these, 3,670 centerlines in the AM peak and 3,669 in the PM peak have valid speeds in both directions and therefore admit directional comparison.
	
	\subsection{Pairing Directions for Asymmetry Detection}
	
	Directional asymmetry is evaluated by comparing speeds in opposite directions along the same physical road segment. Although speeds are observed at the directed centerline level ($skel\_dir$), pairing is performed at the undirected centerline level ($cline\_id$):
	\begin{itemize}
		\item For each $cline\_id$, speeds are pivoted to form paired observations $(speed_{AB}, speed_{BA})$.
		\item Only centerlines with valid speed observations in both directions are retained for asymmetry analysis.
	\end{itemize}
	This procedure ensures that directional comparisons reflect conditions along the same physical infrastructure.
	
	\subsection{Asymmetry Measures and Classification Thresholds}
	
	For each paired centerline, two measures of directional imbalance are computed:
	\begin{itemize}
		\item Speed ratio: $ratio = \min(speed_{AB}/speed_{BA},\, speed_{BA}/speed_{AB}) \in (0,1]$, which captures relative asymmetry in a symmetric manner.
		\item Relative difference: $ratio2 = |speed_{AB} - speed_{BA}| / (speed_{AB} + speed_{BA}) \in [0,1)$, used as a complementary diagnostic.
	\end{itemize}
	
	Directional asymmetry is classified using the following thresholds:
	\begin{itemize}
		\item \textbf{unsym1 (baseline)}: $ratio < 0.5$, corresponding to one direction being at least twice as fast as the other.
		\item \textbf{unsym2 (severe)}: $ratio < 1/3$, corresponding to a speed ratio of at least three-to-one.
		\item \textbf{unsym3 (robustness)}: $ratio2 > 0.6$.
	\end{itemize}
	The baseline analysis uses the unsym1 criterion; the remaining thresholds are used to characterize severity and robustness.
	
	\subsection{Asymmetry Detection Results}
	
	Applying the baseline criterion yields the following peak-specific asymmetry rates:
	\begin{itemize}
		\item AM peak: 405 centerlines (11.0\%) classified as asymmetric; 179 (4.9\%) satisfy the severe asymmetry criterion.
		\item PM peak: 361 centerlines (9.8\%) classified as asymmetric; 145 (4.0\%) satisfy the severe asymmetry criterion.
	\end{itemize}
	
	Figure~\ref{fig:speed_ratio_dist} shows the distribution of speed ratios for AM and PM peaks. Most centerlines have ratios above 0.5, while a lower tail captures segments with pronounced directional imbalance; this tail is thicker in the AM peak than in the PM peak.
	
	\begin{figure}[h!]
		\centering
		\includegraphics[width=0.95\textwidth]{../../../data_work/outputs/manual_centerline_rules_v11/figures/hist_speed_ratio_am_pm.png}
		\caption{Distribution of speed ratios (min(AB/BA, BA/AB)) for AM and PM peak periods. Vertical dashed lines indicate the asymmetry thresholds: $ratio = 0.5$ (unsym1) and $ratio = 1/3$ (unsym2).}
		\label{fig:speed_ratio_dist}
	\end{figure}
	
	\subsection{Temporal Overlap of Asymmetric Segments}
	
	Using the baseline threshold, 638 unique undirected centerlines are classified as asymmetric in at least one peak period:
	\begin{itemize}
		\item Asymmetric in AM only: 277 centerlines (43.4\%).
		\item Asymmetric in PM only: 233 centerlines (36.5\%).
		\item Asymmetric in both AM and PM: 128 centerlines (20.1\%).
	\end{itemize}
	Thus, over 80\% of asymmetric segments appear in only one peak period, indicating strong temporal specificity in directional imbalance.
	
	\subsection{Tidal Lane Candidate Identification}
	
	Tidal lane candidates are defined using two criteria:
	\begin{itemize}
		\item The faster direction reverses between AM and PM peaks.
		\item Both peak periods satisfy the baseline asymmetry criterion ($ratio_{AM} < 0.5$ and $ratio_{PM} < 0.5$).
	\end{itemize}
	Among the 3,668 centerlines with valid bidirectional data in both peaks, 1,032 (28.1\%) exhibit a reversal in the faster direction between peaks, and 14 centerlines (0.4\%) satisfy both conditions and are classified as tidal lane candidates.
	
	\begin{figure}[h!]
		\centering
		\includegraphics[width=0.85\textwidth]{../../../data_work/outputs/manual_centerline_rules_v11/figures/map_tidal_candidates_150km.png}
		\caption{Spatial distribution of tidal lane candidates within 150km of Tiananmen Square.}
		\label{fig:tidal_candidates}
	\end{figure}
	
	\subsection{Spatial Distribution of Asymmetric Segments}
	
	Figures~\ref{fig:asym_am} and~\ref{fig:asym_pm} visualize asymmetric segments during the AM and PM peaks using the baseline threshold. Red segments indicate the AB direction is faster, while blue segments indicate the BA direction is faster.
	
	\begin{figure}[h!]
		\centering
		\begin{minipage}[t]{0.49\textwidth}
			\centering
			\includegraphics[width=\textwidth]{../../../data_work/outputs/manual_centerline_rules_v11/figures/map_asym_AM_unsym1_150km.png}
			\caption{Asymmetric segments in the AM peak ($ratio < 0.5$).}
			\label{fig:asym_am}
		\end{minipage}
		\hfill
		\begin{minipage}[t]{0.49\textwidth}
			\centering
			\includegraphics[width=\textwidth]{../../../data_work/outputs/manual_centerline_rules_v11/figures/map_asym_PM_unsym1_150km.png}
			\caption{Asymmetric segments in the PM peak ($ratio < 0.5$).}
			\label{fig:asym_pm}
		\end{minipage}
	\end{figure}

	\clearpage
%%%%%%%%%%%%%%%%%%%%%%%%%%%%%%%%%%%%%%%%%	

	
	
	
	
	
	
	\section{Grid Construction}
	
	To discretize the urban road network for subsequent aggregation and analysis, we construct three alternative spatial grid systems: square grids, hexagonal grids, and Voronoi (Thiessen) polygons. Each system partitions space using a distinct geometric principle, allowing us to assess the sensitivity of our results to the choice of spatial aggregation units while holding the underlying road network fixed.
	
	\subsection{Square Grid System}
	
	We first construct a regular square grid covering the study area. The resulting system contains 3,357 grid cells with an average area of 8.38~km$^2$.Under this discretization, 64.3\% of directed centerline segments (8,234 out of 12,812) are completely contained within a single grid cell. These segments account for 24.0\% of the total centerline length (4,530.43~km out of 18,884.37~km). 
	
	\subsection{Hexagonal Grid System}
	
	We next construct a hexagonal grid network by Uber's H3 geospatial indexing system, yielding hexagons with an average area of approximately 6.1 km², comparable in scale to our square grid system. This choice ensures that differences across grid systems are not driven by spatial resolution. The H3 grid can also be replaced by a fixed-edge-length grid if required.
	
	This procedure yields 4,604 hexagonal cells, the highest cell count among the three systems. However, due to the irregular shape of hexagons near boundaries, only 58.9\% of centerline segments (7,544) are completely contained within a single cell, representing 19.6\% of total length (3,705.48 km).
	
	\subsection{Voronoi Grid System}
	
	The Voronoi grid system provides a spatially adaptive tessellation anchored directly to the structure of the road network. Unlike uniform square or hexagonal grids, Voronoi cells adjust endogenously to local network density, generating smaller cells in dense urban cores and larger cells in sparsely connected peripheral areas.
	
	\paragraph{Construction.}
	Seed points are drawn from the directed centerline network, using network nodes with degree at least two to prioritize road intersections and major junctions. To prevent excessive clustering of generators in dense regions, we impose a minimum inter-seed distance of 1.5~km through a greedy Poisson-like sampling procedure, yielding approximately 1,200 initial seeds.
	
	To control cell size in sparsely connected regions, we apply an iterative densification procedure. After constructing an initial Voronoi network, cells with area exceeding $A_{\max}=200$~km$^2$ are identified. Within each such cell, an additional seed is placed at the location maximally distant from existing generators, prioritizing network nodes when available. The Voronoi diagram is then recomputed. This process is repeated until all cells satisfy the area constraint, convergence is reached, or a maximum of eight iterations is completed. To limit computational cost, each iteration adds at most 300 new seeds.
	
	This refinement procedure reduces the maximum cell area from over 500~km$^2$ to approximately 195~km$^2$. The final Voronoi grid consists of 1,244 cells with an average area of 22.62~km$^2$, the largest among the three systems. 60.7\% of centerline segments (7,774) are fully contained within a single Voronoi cell, but accounting for only 17.4\% of the total network length (3,284.26~km).
	
	To handle boundary effects, we add auxiliary generators outside the study area and clip the Voronoi tessellation to the analysis boundary, ensuring finite cells.

	\subsection{Grid Comparison}
	
	Table~\ref{tab:grid_comparison} summarizes the key characteristics of the three grid systems. 
	
	Figure~\ref{fig:grid_comparison} provides a visual comparison of the three grid systems overlaid on the directed centerline network, illustrating their differing spatial structures and boundary interactions.
	
	\begin{table}[h!]
		\centering
		\caption{Comparison of three spatial grid systems}
		\label{tab:grid_comparison}
		\begin{tabular}{lcccc}
			\toprule
			Grid Type & Cell Count & Avg. Area (km²) & Segments Within & Length Within \\
			\midrule
			Square & 3,357 & 8.38 & 8,234 (64.3\%) & 4,530.43 km (24.0\%) \\
			Hexagonal & 4,604 & 6.10 & 7,544 (58.9\%) & 3,705.48 km (19.6\%) \\
			Voronoi & 1,244 & 22.62 & 7,774 (60.7\%) & 3,284.26 km (17.4\%) \\
			\bottomrule
		\end{tabular}
		\vspace{0.1cm}
		\footnotesize
		Note: ``Segments Within'' refers to the count and percentage of centerline segments that are completely contained within a single grid cell. ``Length Within'' refers to the total length and percentage of centerline segments that are completely contained within a single grid cell.
	\end{table}
	
	\begin{figure}[h!]
		\centering
		\IfFileExists{../../../00archive/raw_data/gis/map/grid_comparison.png}{
			\includegraphics[width=\textwidth]{../../../00archive/raw_data/gis/map/grid_comparison.png}
		}{
			\fbox{\parbox{0.95\textwidth}{\centering Missing figure file: \texttt{../../../00archive/raw\_data/gis/map/grid\_comparison.png}}}
		}
		\caption{Comparison of square, hexagonal, and Voronoi grid systems}
		\label{fig:grid_comparison}
	\end{figure}


	%%%%%%%%%%%%%%%%%%%%%%%%%%%%%%%%%%%%%%%%%	
	%  Model =========================================================

	\clearpage
	\section{Model}
	
	\subsection{Setup}
	This model is a urban commuting model with directed travel time. The city consists of a finite set of locations (grid cells) $\mathcal N=\{1,\dots,N\}$ connected by a directed graph $\mathcal E\subseteq\mathcal N\times\mathcal N$. For each directed link $(k,l)\in\mathcal E$, let $\ell_{kl}>0$ denote its physical length, $v^0_{kl}>0$ its free-flow (non-congested) speed, and $t_{kl}>0$ the realized equilibrium travel time. Directionality allows $\ell_{kl}\neq\ell_{lk}$, $v^0_{kl}\neq v^0_{lk}$, and $t_{kl}\neq t_{lk}$ in general.
	
	An aggregate labor endowment $\bar L$ is allocated across residential locations $\{L_i^R\}_{i\in\mathcal N}$ and workplace locations $\{L_i^F\}_{i\in\mathcal N}$, with $\sum_i L_i^R=\sum_i L_i^F=\bar L$. The model is parameterized by $\theta>0$, governing heterogeneity in route choice; spillover elasticities $\alpha$ and $\beta$; and a traffic congestion elasticity $\lambda\geq 0$.
	
	Each location $i$ is endowed with baseline productivity $\bar A_i>0$ and baseline amenity $\bar u_i>0$. Local productivity and amenities depend on endogenous activity according to
	\begin{equation}
		A_i=\bar A_i (L_i^F)^{\alpha}, \qquad u_i=\bar u_i (L_i^R)^{\beta}.
		\label{eq:fundamentals}
	\end{equation}
	
	\subsection{Individuals: residence, workplace, and route choice}
	An individual chooses a residence $i\in\mathcal N$, a workplace $j\in\mathcal N$, and a commuting route $r\in\mathcal R_{ij}$, where a route is a finite ordered sequence of nodes $r=(r_0=i,r_1,\dots,r_K=j)$. The payoff associated with choice $(i,j,r)$ is
	\begin{equation}
		V_{ij,r}=\left(\frac{u_i w_j}{\prod_{m=1}^{K} t_{r_{m-1},r_m}}\right)\varepsilon_{ij,r},
		\label{eq:utility}
	\end{equation}
	where $w_j$ is the wage at workplace $j$ and $\varepsilon_{ij,r}$ is i.i.d.\ Fr\'echet with shape parameter $\theta$.
	
	\paragraph{Choice probabilities.}
	Given Fr\'echet heterogeneity, the probability that an individual chooses $(i,j,r)$ is
	\begin{equation}
		\pi_{ij,r}
		=
		\frac{\left(u_i w_j \prod_{m=1}^{K} t_{r_{m-1},r_m}^{-1}\right)^{\theta}}
		{\sum_{i',j'}\sum_{r'\in\mathcal R_{i'j'}}
			\left(u_{i'} w_{j'} \prod_{m=1}^{K'} t_{r'_{m-1},r'_m}^{-1}\right)^{\theta}}.
		\label{eq:routeprob}
	\end{equation}
	
	\subsection{Aggregation over routes and commuting costs}
	Let $\mathcal R_{ij}$ denote the (countably infinite) set of feasible routes connecting residence $i$ to workplace $j$. Aggregating \eqref{eq:routeprob} over all routes and multiplying by aggregate population $\bar L$ yields the commuting flow
	\begin{equation}
		L_{ij}=\bar L \sum_{r\in\mathcal R_{ij}} \pi_{ij,r}.
		\label{eq:flowdef}
	\end{equation}
	Since $u_i^\theta w_j^\theta$ does not vary across routes, define the bilateral commuting cost index
	\begin{equation}
		\tau_{ij}^{-\theta}
		\equiv
		\sum_{r\in\mathcal R_{ij}}
		\prod_{m=1}^{K} t_{r_{m-1},r_m}^{-\theta},
		\label{eq:tau_def}
	\end{equation}
	which summarizes all feasible routes between $i$ and $j$.
	
	\subsection{Network representation of commuting costs}
	Define the weighted adjacency matrix $A=[a_{kl}]$ with $a_{kl}=t_{kl}^{-\theta}$. For any $K\ge0$, $(A^K)_{ij}$ collects the total weight of all routes of length $K$ from $i$ to $j$. If the spectral radius of $A$ is less than one,
	\begin{equation}
		\tau_{ij}^{-\theta}=\sum_{K=0}^{\infty}(A^K)_{ij}=[(I-A)^{-1}]_{ij}\equiv b_{ij},
		\label{eq:tau}
	\end{equation}
	so that $\tau_{ij}=b_{ij}^{-1/\theta}$. Directionality of the network implies $\tau_{ij}\neq\tau_{ji}$ in general.
	
	\subsection{Firms}
	Firms operate under perfect competition at each location $j$ and use labor as the sole input. Wages equal local productivity,
	\begin{equation}
		w_j=A_j=\bar A_j (L_j^F)^{\alpha}.
		\label{eq:wage}
	\end{equation}
	
	\subsection{Commuting flows and market clearing}
	Using \eqref{eq:tau}, commuting flows take the gravity form
	\begin{equation}
		L_{ij}
		=
		\tau_{ij}^{-\theta} u_i^\theta w_j^\theta
		\frac{\bar L}{\bar W^\theta},
		\label{eq:gravity}
	\end{equation}
	where
	\begin{equation}
		\bar W
		=
		\left(
		\sum_{i,j}\tau_{ij}^{-\theta}u_i^\theta w_j^\theta
		\right)^{1/\theta}
		\label{eq:welfare}
	\end{equation}
	is expected equilibrium welfare. Residential and employment populations satisfy
	\begin{equation}
		L_i^R=\sum_j L_{ij}, \qquad L_j^F=\sum_i L_{ij}.
		\label{eq:marketclearing}
	\end{equation}
	
	\subsection{Transportation Network, Lane Capacity, and Congestion}

	\paragraph{Free-flow travel time.}
	For each directed link $(k,l)\in\mathcal E$, the non-congested travel time is fully determined by its physical characteristics:
	\begin{equation}
		\bar{t}_{kl} = \frac{\ell_{kl}}{v^0_{kl}}.
		\label{eq:tbar}
	\end{equation}
	Both $\ell_{kl}$ and $v^0_{kl}$ may differ by direction, so in general $\bar{t}_{kl}\neq\bar{t}_{lk}$.

	\paragraph{Traffic on a directed link.}
	Let $\pi^{kl}_{ij}$ denote the share of commuters from $i$ to $j$ whose optimal route passes through directed link $(k,l)$. Directed traffic on link $(k,l)$ is
	\begin{equation}
		\phi_{kl}
		=
		\sum_{i\in\mathcal N}\sum_{j\in\mathcal N}
		L_{ij}\,\pi^{kl}_{ij}.
		\label{eq:traffic_kl}
	\end{equation}

	\paragraph{Congestion technology.}
	Following the spirit of Allen \& Arkolakis (2022), realized travel time is increasing in equilibrium traffic per lane:
	\begin{equation}
		t_{kl}
		=
		\bar{t}_{kl}
		\left(
		\frac{\phi_{kl}}{n_{kl}}
		\right)^{\!\lambda},
		\qquad \lambda \geq 0,
		\label{eq:tkl_cong}
	\end{equation}
	where $n_{kl}$ is the directed lane count in direction $(k\to l)$. Substituting \eqref{eq:tbar}, this expands to
	\begin{equation}
		t_{kl}
		=
		\frac{\ell_{kl}}{v^0_{kl}}
		\left(
		\frac{\phi_{kl}}{n_{kl}}
		\right)^{\!\lambda}.
		\label{eq:tkl_expand}
	\end{equation}
	When $\lambda=0$, $t_{kl}=\bar{t}_{kl}$ and the model reduces to the frictionless case. When $\lambda>0$, higher traffic per lane raises travel times, creating a feedback between commuting decisions and congestion that is resolved jointly in equilibrium.

	\paragraph{Directed lanes and tidal asymmetry.}
	We define $n_{kl}$ as the directed lane count in direction $(k\to l)$, which may differ from $n_{lk}$. Under a standard symmetric allocation $n_{kl}\approx n_{lk}$. A tidal lane configuration redistributes capacity across directions so that $n_{kl}\neq n_{lk}$. This induces asymmetric congestion through \eqref{eq:tkl_cong} and hence asymmetric bilateral commuting costs $\tau_{ij}\neq\tau_{ji}$ through \eqref{eq:tau}.
	
	\subsection{Equilibrium}

	An equilibrium is a collection
	$\bigl(\{L_i^R,L_i^F\}_{i\in\mathcal N},\,\{L_{ij}\}_{i,j\in\mathcal N},\,\{\tau_{ij}\}_{i,j\in\mathcal N},\,\{t_{kl}\}_{(k,l)\in\mathcal E},\,\{\phi_{kl}\}_{(k,l)\in\mathcal E}\bigr)$
	such that:
	\begin{enumerate}[label=(\roman*)]
		\item given directed travel times $\{t_{kl}\}$, bilateral commuting costs $\{\tau_{ij}\}$ are generated by optimal routing on the directed network as in \eqref{eq:tau};
		\item given $\{\tau_{ij}\}$, commuting flows $\{L_{ij}\}$ satisfy the gravity equation \eqref{eq:gravity};
		\item residential and employment populations clear markets as in \eqref{eq:marketclearing};
		\item directed traffic $\{\phi_{kl}\}$ is induced by equilibrium commuting flows and route choice as in \eqref{eq:traffic_kl}; and
		\item realized travel times $\{t_{kl}\}$ satisfy the congestion technology \eqref{eq:tkl_cong}.
	\end{enumerate}
	Conditions (i)--(iii) are standard to quantitative spatial commuting models; conditions (iv)--(v) follow Allen \& Arkolakis (2022). The bilateral commuting cost $\tau_{ij}$ is \emph{endogenous}: it depends on the full network of link costs $\{t_{kl}\}$, which are shaped by the distribution of commuting activity through congestion. When $\lambda=0$, conditions (iv)--(v) are vacuous and the model reduces to the standard no-congestion case.

	Let $l_i^R=L_i^R/\bar L$ and $l_i^F=L_i^F/\bar L$ denote population shares. Given equilibrium transportation costs $\{\tau_{ij}\}$, the market-clearing conditions yield the equilibrium allocation $(l^R,l^F)$:
	\begin{align}
		(l_i^R)^{1-\theta\beta} &= \chi \sum_j \tau_{ij}^{-\theta}\,\bar u_i^\theta \bar A_j^\theta (l_j^F)^{\alpha\theta}, \label{eq:eq_R}\\
		(l_i^F)^{1-\theta\alpha} &= \chi \sum_j \tau_{ji}^{-\theta}\,\bar u_j^\theta \bar A_i^\theta (l_j^R)^{\beta\theta}, \label{eq:eq_F}
	\end{align}
	where $\chi\equiv \bar L^{\theta(\alpha+\beta)}/\bar W^{\theta}$. Together with conditions (i)--(v), equations \eqref{eq:eq_R}--\eqref{eq:eq_F} characterize the full equilibrium. Route choice, congestion, and commuting flows are jointly determined: commuting patterns determine traffic $\phi_{kl}$, traffic determines travel times $t_{kl}$ through \eqref{eq:tkl_cong}, and travel times determine bilateral commuting costs $\tau_{ij}$ through \eqref{eq:tau}.





	%%%%%%%%%%%%%%%%%%%%%%%%%%%%%%%%%%%%%%%%%
	
	\clearpage
	\section{Model Estimation}
	
	\subsection{Data}
	\paragraph{Spatial units and network nodes.} The set of locations is a finite collection of nodes $\mathcal N=\{1,\dots,N\}$, where each node corresponds to a spatial cell. We implement the model on three alternative grid systems (square, hexagon, and Voronoi-type partitions). For each grid system, a node $i\in\mathcal N$ is identified by a unique \texttt{grid\_id}. The empirical procedures below are identical across grid systems; the choice of grid mainly changes the discretization of the city and the induced network topology.
	
	\paragraph{Directed edges, link costs, and lane counts.}
	For each grid system, we construct a directed adjacency set $\mathcal E\subseteq\mathcal N\times\mathcal N$. A directed edge $(k,l)\in\mathcal E$ represents an admissible move from node $k$ to node $l$ along the observed road infrastructure. For each directed grid link $(k,l)$ we construct three objects from the aggregated centerline data:\footnote{When a directed road segment is entirely contained within a single grid cell, it contributes only to within-grid stocks (e.g., lane-km and average speed). If a segment intersects multiple grid cells, it is first split at grid boundaries into sub-segments, each uniquely assigned to one grid. These sub-segments are then ordered along the original geometry to recover the sequence of traversed grids. Directed links are constructed between consecutive grids (e.g., $A \to B$, $B \to C$), with link attributes computed from the corresponding sub-segments, including length, lane-km, and average speed. Aggregating across all segments yields a grid-level directed network characterized by bilateral travel costs, speeds, and lane counts.}
	\begin{itemize}
		\item \textbf{Grid distance} $\text{dist}_{kl}$: centroid-to-centroid distance between adjacent grids.
		\item \textbf{Observed (congestion-inclusive) speed} $v_{kl}^{\text{obs}}$: harmonic mean AM speed across all centerline pieces crossing from grid $k$ to grid $l$.
		\item \textbf{Free-flow speed} $v_{kl}^{0}$: harmonic mean speed constructed from raw observations in the `22:00--05:00' window, aggregated first to directed centerlines and then to grid links.
		\item \textbf{Directed lane count} $\text{lanes}_{kl}$: length-weighted average of approximate lane counts derived from the road type attribute (\texttt{roadtype}) of matched centerline segments, using a standard road-type-to-lane lookup table.
	\end{itemize}
	The \emph{observed} (congestion-inclusive) link travel time is then
	\begin{equation}\label{eq:data_tkl}
		t_{kl}\;\equiv\; \frac{\text{dist}_{kl}}{v_{kl}^{\text{obs}}},
	\end{equation}
	measured in minutes and reflecting equilibrium congestion. In the current implementation, the free-flow travel time is observed from the nighttime speed window:
	\begin{equation}
		t_{kl}^{ff}\;\equiv\; \frac{\text{dist}_{kl}}{v_{kl}^{0}}.
	\end{equation}
	Directionality is preserved throughout: $t_{kl}$, $t_{lk}$, $t_{kl}^{ff}$, $t_{lk}^{ff}$, $n_{kl}$, and $n_{lk}$ can all differ, directly capturing tidal asymmetry in road use and lane capacity.
	
	Figure \ref{fig:gridlink_example} illustrates the resulting grid-level representation for the 3km square grid in the morning peak (AM): the left panel shows observed centerline speeds, and the right panel maps them into directed grid-to-grid link travel times. We construct analogous link measures for both AM and PM periods and for all grid systems.
	
	\begin{figure}[h!]
		\centering
		\includegraphics[width=0.95\textwidth]{../../../data_work/outputs/raw_rebuild_validation/figures/grid_links_square3km_AM.png}
		\caption{Example of grid-level directed link construction and implied travel times: 3km square grid, AM peak.}
		\label{fig:gridlink_example}
	\end{figure}


	
	\paragraph{Commuting flows and marginal populations.} Using the commuting matrix from Baidu Map, we assign each home coordinate and work coordinate to a grid cell via point-in-polygon matching. We then aggregate to obtain grid-to-grid commuting flows $L_{ij}$:
	\begin{equation}\label{eq:data_Lij}
		L_{ij}\equiv \sum_{\text{commuters }(i\to j)} \text{pop},
	\end{equation}
	where \text{pop} is the number of commuters associated with each home--work pair. Residential and workplace marginal totals are computed as
	\begin{equation}\label{eq:data_margins}
		L_i^R=\sum_{j\in\mathcal N}L_{ij},\qquad L_j^F=\sum_{i\in\mathcal N}L_{ij},\qquad \bar L=\sum_{i\in\mathcal N}L_i^R=\sum_{j\in\mathcal N}L_j^F.
	\end{equation}
	When available, we also retain commute-mode counts (walk, bike, subway, bus, car) to describe coverage and heterogeneity in observed mode choice, noting that the sum of identified modes may be below total \text{pop}.
	
	\paragraph{Reachability filter.} Because not all node pairs $(i,j)$ with commuting flow data are mutually reachable, we restrict attention to the largest weakly connected component before computing bilateral commuting costs.\footnote{In the square 3km grid system, the raw commuting data contain 442,796 distinct grid--grid origin--destination pairs, covering a total commuting population of 7.96 million individuals. After restricting attention to grid pairs whose home and work locations lie in the same weakly connected component of the directed road network, 374,850 OD pairs remain (84.7\% of all pairs), accounting for 7.59 million commuters (95.3\% of the total population). This step therefore primarily removes low-volume and peripheral commuting pairs while preserving the vast majority of economically relevant commuting flows.} This step ensures that bilateral commuting costs are well-defined on the analysis sample and avoids numerical artifacts driven by disconnected nodes.

	\subsection{Parameters}
	
	\paragraph{Calibrated parameters.}
	The baseline implementation uses external calibration for $(\theta,\alpha,\beta)$.
	The parameter $\theta>0$ governs Fr\'echet heterogeneity in route choice, while
	$\alpha$ and $\beta$ capture productivity and amenity spillovers across locations,
	respectively.

	As these parameters are standard in the quantitative urban and commuting literature, we adopt benchmark values from existing studies. In particular, we follow the estimates in Ahlfeldt et al.\ (2015), which are obtained in an urban spatial equilibrium model with endogenous commuting and land markets.
	Specifically, we set $\theta = 6.83$, $\alpha = -0.12$, and $\beta = -0.10$.

	These values lie within the range commonly used in the urban economics literature and imply economically meaningful but moderate agglomeration and amenity externalities. In the current codebase, a two-way fixed-effects gravity regression is also run as a diagnostic for $\theta$, but the baseline implementation does not feed that estimate back into the equilibrium solver.

	\paragraph{Congestion elasticity $\lambda$.}
	The congestion elasticity $\lambda\geq 0$ governs how strongly realized travel time responds to traffic per lane, and thus determines the magnitude of the travel-time improvement from adding a lane. In the current implementation, we estimate a reduced-form cross-sectional proxy based on the free-flow and observed grid-edge times:
	\begin{equation}\label{eq:est_lambda}
		\ln\!\left(\frac{t_{kl}^{obs}}{t_{kl}^{ff}}\right)
		=
		\eta_{\mathrm{lane\text{-}bin}(kl)}
		+
		\lambda\,\ln\!\left(\frac{\phi_{kl}^{obs}}{n_{kl}}\right)
		+ \varepsilon_{kl},
	\end{equation}
	where $\phi_{kl}^{obs}$ is an observed-flow proxy obtained by assigning the observed OD matrix to shortest paths on the observed directed grid graph. The lane-bin fixed effects absorb coarse capacity differences across links. This estimate should be interpreted as a weak calibration device rather than a fully structural identification of congestion.
	In the current square-grid baseline, this regression yields $\hat\lambda=0.01$, which is the imposed lower bound in the calibration routine; we therefore treat the resulting congestion response as conservative.

	\paragraph{Robustness.} To assess sensitivity, we consider alternative calibrations within plausible ranges: (i) varying $\theta$ (e.g.\ low/medium/high elasticities) to evaluate how strongly results depend on route-choice dispersion; (ii) varying $\alpha$ and $\beta$ within literature-reported bounds to gauge the role of agglomeration and residential spillovers; (iii) varying $\lambda$ around the estimated baseline to assess the sensitivity of counterfactual welfare gains to the strength of traffic congestion.
	
	\paragraph{Theoretically, we can estimate $\theta$ via a gravity equation with OD fixed effects}
	Although we externally calibrate $\theta$ in the baseline, one can estimate $\theta$ from commuting flows using a gravity regression with origin and destination fixed effects. The gravity equation implied by the model is
	\begin{equation}\label{eq:est_gravity}
		L_{ij}=\tau_{ij}^{-\theta}u_i^\theta w_j^\theta\frac{\bar L}{\bar W^\theta},
	\end{equation}
	where $\tau_{ij}$ is the bilateral commuting cost index, and $\bar W$ is expected welfare. Taking logs and absorbing all origin and destination terms into fixed effects yields
	\begin{equation}\label{eq:est_gravity_fe}
		\log L_{ij}=-\theta \log \tau_{ij}+\mu_i+\nu_j+\varepsilon_{ij},
	\end{equation}
	where $\mu_i$ and $\nu_j$ capture the composite effects of amenities, productivity, and endogenous scale terms. Because $\tau_{ij}$ depends on $\theta$ through the network aggregator (see equation \eqref{eq:tau}), a practical implementation uses an outer grid search: for each candidate $\theta$, compute $\tau_{ij}(\theta)$ on the OD sample and estimate \eqref{eq:est_gravity_fe}, selecting $\theta$ that optimizes fit.
	
	\subsection{Inversion}
	
	\paragraph{Inverting fundamentals $(\bar u_i,\bar A_i)$ given calibrated parameters}
	Given $(\theta,\alpha,\beta)$, observed $(L^R,L^F)$ from \eqref{eq:data_margins}, and network-implied $\{\tau_{ij}\}$, we invert for baseline fundamentals $(\bar u_i,\bar A_i)$ consistent with equilibrium. The equilibrium system can be written as two sets of node-level conditions (in population shares $l_i^R=L_i^R/\bar L$ and $l_i^F=L_i^F/\bar L$):
	\begin{align}
		(l_i^R)^{1-\theta\beta} &= \chi \sum_{j\in\mathcal N}\tau_{ij}^{-\theta}\bar u_i^\theta \bar A_j^\theta (l_j^F)^{\alpha\theta},\label{eq:invert_R}\\
		(l_i^F)^{1-\theta\alpha} &= \chi \sum_{j\in\mathcal N}\tau_{ji}^{-\theta}\bar u_j^\theta \bar A_i^\theta (l_j^R)^{\beta\theta},\label{eq:invert_F}
	\end{align}
	where $\chi>0$ is a common scalar proportional to equilibrium welfare (and pinned down by a normalization). 
	
		\subsection{Counterfactual: Tidal-Lane Policies as Directed Lane Reallocations}

	We model a tidal lane policy as a reallocation of directed lane capacity on a targeted set of links $\mathcal{E}^\star\subseteq\mathcal{E}$, identified from the reduced-form evidence in Section~\ref{sec:asym_sec}. The policy assigns $\Delta_{kl}\geq 0$ lanes from the off-peak direction $(l\to k)$ to the peak-demand direction $(k\to l)$:
	\begin{equation}\label{eq:cf_lanes}
		n_{kl}^{\mathrm{cf}} = n_{kl} + \Delta_{kl},
		\qquad
		n_{lk}^{\mathrm{cf}} = n_{lk} - \Delta_{kl},
		\qquad (k,l)\in\mathcal{E}^\star,
	\end{equation}
	with all other links unchanged. This is a lane-neutral reallocation: total road capacity $n_{kl}+n_{lk}$ is preserved, but the directional distribution shifts in favor of the peak-demand direction.

	The policy enters the model directly through the congestion technology \eqref{eq:tkl_cong}: adding lanes to direction $(k\to l)$ reduces traffic per lane $\phi_{kl}/n_{kl}$, lowering travel time $t_{kl}$; simultaneously, removing lanes from the reverse direction $(l\to k)$ raises $\phi_{lk}/n_{lk}$, increasing $t_{lk}$. The magnitude of both effects is governed by $\lambda$. Because $\phi_{kl}$ is itself endogenous, commuters re-route in response to the changed travel times, and the two effects interact throughout the full network.

	Given the counterfactual lane counts $\{n_{kl}^{\mathrm{cf}}\}$, we solve the counterfactual equilibrium (conditions (i)--(v) with $n_{kl}^{\mathrm{cf}}$ in place of $n_{kl}$), obtain $\tau_{ij}^{\mathrm{cf}}$ from \eqref{eq:tau}, and solve the equilibrium system \eqref{eq:eq_R}--\eqref{eq:eq_F} holding calibrated parameters $(\theta,\alpha,\beta,\lambda)$ and inverted fundamentals $(\bar{u}_i,\bar{A}_i)$ fixed. Comparing baseline and counterfactual outcomes delivers changes in directional commuting costs, spatial allocations $(L^{R,\mathrm{cf}},L^{F,\mathrm{cf}})$, and aggregate welfare $\hat{\bar{W}}$.

	The targeted set $\mathcal{E}^\star$ is constructed by mapping road centerlines with detected tidal patterns (Section~\ref{sec:asym_sec}) to their corresponding directed links in the grid-level network.
	
	\subsection{Data and variables}
	Table~\ref{tab:data_mapping} summarizes the correspondence between model variables
	and the underlying data constructed in the grid-level preprocessing pipeline.
	All grid-based objects are constructed for three alternative spatial systems
	(square grid, hexagonal grid, and administrative grid), using the same procedures.
	
	\begin{table}[htbp]\centering
		\caption{Model variables and data construction}
		\label{tab:data_mapping}
		\begin{tabular}{lll}
			\hline
			Variables & Symbol & Data construction \\
			\hline
			Locations (nodes) 
			& $i\in\mathcal N$ 
			& Grid centroids (square / hex / Voronoi) \\
			
			Directed links 
			& $(k,l)\in\mathcal E$ 
			& Adjacency of neighboring grid cells \\
			
			Directed travel time 
			& $t_{kl}$ 
			& Grid centroid distance divided by harmonic mean crossing speed \\

			Free-flow travel time
			& $t_{kl}^{ff}$
			& Grid centroid distance divided by nighttime (`22:00--05:00') harmonic speed \\
			
			Network weights 
			& $A_{kl}=t_{kl}^{-\theta}$ 
			& Constructed from $t_{kl}$ given $\theta$ \\
			
			OD commuting cost 
			& $\tau_{ij}$ 
			& Shortest-path travel time on the directed grid graph \\
			
			Residential population 
			& $L_i^{R}$ 
			&  population counts by grid \\
			
			Employment population 
			& $L_i^{F}$ 
			&  workers counts by grid \\
			
			Baseline fundamentals 
			& $(\bar u_i,\bar A_i)$ 
			& Inverted from equilibrium conditions \\
			
			Counterfactual objects 
			& $(t_{kl}^{\mathrm{cf}},\tau_{ij}^{\mathrm{cf}})$ 
			& Directional shocks to selected links \\
			\hline
		\end{tabular}
	\end{table}
	
	Figure~\ref{fig:grid_link_example} illustrates how observed directional road-level speeds are aggregated into grid-to-grid link travel costs in the square 3km grid during the morning peak. The left panel plots observed centerline speeds, while the right panel shows the resulting grid-level directed network, where link travel times are computed from centroid distance divided by the harmonic mean of crossing-piece speeds. The red outline highlights the outer boundary of the largest weakly connected component of the grid-level road network, which defines the set of grids retained for subsequent commuting-cost computations and excludes peripheral grids that are not mutually reachable through the observed road infrastructure.
	
	\begin{figure}[h!]
		\centering
		\includegraphics[width=1\textwidth]{../../../data_work/outputs/raw_rebuild_validation/figures/aa_style_square3km_AM_2panel.png}
		\caption{From road segments to grid-level link costs: square 3km grid, AM peak.
			Left panel plots observed centerline speeds. Right panel shows the corresponding
			grid-to-grid network, with directed link travel times constructed from aggregated
			segment-level information.}
		\label{fig:grid_link_example}
	\end{figure}
	
	Figure~\ref{fig:grid_speed_ratio} documents the extent of directional speed asymmetry
	within the grid system. The figure plots the distribution of the speed ratio
	$\min(v_{ab}/v_{ba},\, v_{ba}/v_{ab})$ across grid pairs during the AM and PM peak periods.
	The resulting distributions closely mirror those obtained at the centerline level
	(Figure~\ref{fig:speed_ratio_dist}), indicating that the grid aggregation preserves
	the key patterns of directional congestion observed in the underlying road network.
	
	\begin{figure}[h!]
		\centering
		\begin{minipage}[t]{0.49\textwidth}
			\centering
			\includegraphics[width=\textwidth]{../../../data_work/outputs/raw_rebuild_validation/figures/asym_hist_AM_square3km.png}
		\end{minipage}
		\hfill
		\begin{minipage}[t]{0.49\textwidth}
			\centering
			\includegraphics[width=\textwidth]{../../../data_work/outputs/raw_rebuild_validation/figures/asym_hist_PM_square3km.png}
		\end{minipage}
		\caption{Distribution of directional speed asymmetry at the grid level.
			The left panel reports the AM distribution and the right panel reports the PM distribution of
			$\min(v_{ab}/v_{ba},\, v_{ba}/v_{ab})$ for directed grid-to-grid links in the square 3km system.}
		\label{fig:grid_speed_ratio}
	\end{figure}

	\clearpage 
	\section{Conclusion}
	\begin{itemize}
		\item Contribution: incorporate directional, grid-level networks into a quantitative spatial framework.
		\item Policy relevance: welfare analysis of tidal lanes in Chinese megacities.
	\end{itemize}

\end{document}
